\section{MCMC --- the MMSE estimator}
\label{sec:mcmc}

\paragraph{How it works.}
MCMC is the only \emph{MMSE} method here: it returns the posterior mean $\mathbb{E}[f\mid g]$, not
a mode, by a Metropolis-Hastings chain that it averages after a burn-in:
\begin{quote}\small
$z \leftarrow f + (H^{\!\top}\!H+\lambda_{rr} I)^{-1} H^{\!\top}(g-Hf)$ (pull to the data);\quad
$z \leftarrow z + \sigma\varepsilon$ (explore);\quad $z \leftarrow D_\sigma(\max(z,0))$ (denoise);\\
accept $z$ with probability $\min(1, e^{-\Delta D/\beta})$; average the states after burn-in.
\end{quote}
The prior enters through the denoiser; the acceptance is a selector and $\beta$ sets how strict it
is (the target is a \emph{selective} rate ${\approx}0.25$--$0.6$, not saturation). It has no dual
variable --- like PnP-HQS, and unlike ADMM-PnP.

\paragraph{Parameters.}
The denoiser and $\sigma$ (exploration \emph{and} denoising strength, a fraction of the peak, so
it transfers between problems). But the decisive knob is $\lambda_{rr}$, the \textbf{data-step
ridge}: it corrects the directions with $s_i^2 \gg \lambda_{rr}$ toward the data and leaves the
rest. Its useful window is $[s_3^2,\,s_2^2] \approx [0.3,\,36]$. The chain also needs enough
samples (\texttt{max\_iter} with a 25\% burn-in) and a calibrated temperature.

\paragraph{Standard use.}
The out-of-box default (TV~Bregman, $\sigma=0.02$, automatic $\lambda_{rr}=s_1=38.9$,
\texttt{max\_iter}=300) is \emph{weak}: $\lambda_{rr}=s_1$ sits \emph{above} the useful window, so
the data-step barely pulls and the chain drifts.

\paragraph{Optimal use.}
Bring $\lambda_{rr}$ \emph{into} the window --- $\lambda_{rr}\approx5$ --- and take more samples.
This is the single largest lever (it moved the reconstruction from ${\approx}$ the drifting default
to well below the ridge start). The best denoiser is again object-dependent: \textbf{Gaussian} at
low noise, shifting to \textbf{TV~Bregman} (membrane) or \textbf{DCT} (filaments, point-like) under
noise. The measured figures per structure:
\IfFileExists{tables/algo_mcmc.tex}{\begin{center}\begin{minipage}{\linewidth}\centering
\captionof{table}{MCMC on MA-TIRF --- reconstruction metrics per ground truth and noise level.}
\begin{adjustbox}{max width=\linewidth}
\begin{tabular}{llllllll}
\toprule
dataset & noise & params & NMSE & PSNR & Depth Error (nm) & Stack Recovery & chi2\_ratio \\
\midrule
cell & 0 & denoiser=TV Bregman, sigma=0.02, lambda\_rr=38.863, iter=300 & 0.428 & 17.5 & 20 & 1 & 3.22e-05 \\
cell & 0 & denoiser=Gaussian, sigma=0.01, lambda\_rr=5, iter=300 & 0.227 & 20.2 & 11 & 1 & 1.35e-05 \\
fibres & 0 & denoiser=TV Bregman, sigma=0.02, lambda\_rr=38.863, iter=300 & 0.314 & 29.1 & 45 & 0.991 & 0.811 \\
fibres & 0 & denoiser=Gaussian, sigma=0.01, lambda\_rr=5, iter=300 & 0.115 & 33.4 & 23 & 0.991 & 0.121 \\
vesicles & 0 & denoiser=TV Bregman, sigma=0.02, lambda\_rr=38.863, iter=300 & 0.425 & 29.8 & 38 & 1 & 0.52 \\
vesicles & 0 & denoiser=Gaussian, sigma=0.01, lambda\_rr=5, iter=300 & 0.166 & 33.9 & 19 & 1 & 0.0914 \\
cell & 0.02 & denoiser=TV Bregman, sigma=0.02, lambda\_rr=38.863, iter=300 & 0.436 & 17.4 & 20 & 0.985 & 0.724 \\
cell & 0.02 & denoiser=TV Bregman, sigma=0.01, lambda\_rr=5, iter=300 & 0.235 & 20.1 & 13 & 0.742 & 0.894 \\
fibres & 0.02 & denoiser=TV Bregman, sigma=0.02, lambda\_rr=38.863, iter=300 & 0.429 & 27.7 & 61 & 0.966 & 0.735 \\
vesicles & 0.02 & denoiser=TV Bregman, sigma=0.02, lambda\_rr=38.863, iter=300 & 0.512 & 29.0 & 61 & 1 & 0.609 \\
cell & 0.05 & denoiser=TV Bregman, sigma=0.02, lambda\_rr=38.863, iter=300 & 0.512 & 16.7 & 26 & 0.924 & 0.64 \\
cell & 0.05 & denoiser=TV Bregman, sigma=0.01, lambda\_rr=5, iter=300 & 0.453 & 17.2 & 31 & 0.561 & 0.822 \\
fibres & 0.05 & denoiser=TV Bregman, sigma=0.02, lambda\_rr=38.863, iter=300 & 0.772 & 25.2 & 96 & 0.793 & 0.477 \\
vesicles & 0.05 & denoiser=TV Bregman, sigma=0.02, lambda\_rr=38.863, iter=300 & 0.799 & 27.1 & 110 & 1 & 0.435 \\
esoubies & native & denoiser=TV Bregman, sigma=0.01, lambda\_rr=5, iter=300 & — & — & — & — & 53.2 \\
\bottomrule
\end{tabular}\end{adjustbox}\end{minipage}\end{center}
\begin{center}\begin{minipage}{\linewidth}\centering
\captionof{table}{MCMC on deconvolution --- reconstruction metrics per noise level.}
\begin{adjustbox}{max width=\linewidth}
\begin{tabular}{lllllll}
\toprule
dataset & noise & params & NMSE & PSNR & SSIM & chi2\_ratio \\
\midrule
img\_001 & 0.02 & denoiser=TV Bregman, sigma=0.03 & 0.0245 & 23.4 & 0.596 & 0.98 \\
\bottomrule
\end{tabular}\end{adjustbox}\end{minipage}\end{center}
}%
  {\emph{[pending the run — \texttt{python -m benchmarks analyze}]}}

\noindent Which denoiser wins on each object (point-like / filaments / membrane):
\IfFileExists{tables/denoisers_mcmc.tex}{\paragraph{matirf}
\begin{center}\begin{adjustbox}{max width=\linewidth}
\begin{tabular}{llllllll}
\toprule
structure & noise & denoiser & NMSE & PSNR & Depth Error (nm) & Stack Recovery & chi2\_ratio \\
\midrule
cell & 0 & \textbf{Gaussian} & 0.227 & 20.2 & 11 & 1 & 1.35e-05 \\
cell & 0 & TV Bregman & 0.238 & 20.0 & 10 & 0.697 & 1.21e-05 \\
cell & 0 & DCT & 0.302 & 19.0 & 14 & 0.0152 & 2.32e-05 \\
fibres & 0 & \textbf{Gaussian} & 0.115 & 33.4 & 23 & 0.991 & 0.121 \\
fibres & 0 & TV Bregman & 0.213 & 30.8 & 33 & 0.905 & 0.44 \\
fibres & 0 & DCT & 0.282 & 29.5 & 57 & 0 & 0.441 \\
vesicles & 0 & \textbf{Gaussian} & 0.166 & 33.9 & 19 & 1 & 0.0914 \\
vesicles & 0 & TV Bregman & 0.305 & 31.3 & 23 & 0.5 & 0.23 \\
vesicles & 0 & DCT & 0.45 & 29.6 & 62 & 0 & 0.214 \\
cell & 0.02 & \textbf{TV Bregman} & 0.235 & 20.1 & 13 & 0.742 & 0.894 \\
cell & 0.02 & DCT & 0.317 & 18.8 & 17 & 0.0455 & 0.925 \\
cell & 0.02 & Gaussian & 0.402 & 17.7 & 28 & 0.788 & 1.2 \\
fibres & 0.02 & \textbf{DCT} & 0.511 & 27.0 & 81 & 0 & 0.897 \\
fibres & 0.02 & TV Bregman & 0.589 & 26.3 & 67 & 0.612 & 0.954 \\
fibres & 0.02 & Gaussian & 0.791 & 25.1 & 61 & 0.534 & 0.986 \\
vesicles & 0.02 & \textbf{DCT} & 0.588 & 28.4 & 87 & 0 & 0.736 \\
vesicles & 0.02 & TV Bregman & 0.633 & 28.1 & 78 & 0.5 & 0.822 \\
vesicles & 0.02 & Gaussian & 0.797 & 27.1 & 74 & 0.5 & 0.922 \\
cell & 0.05 & \textbf{TV Bregman} & 0.453 & 17.2 & 31 & 0.561 & 0.822 \\
cell & 0.05 & DCT & 0.641 & 15.7 & 57 & 0.0455 & 0.875 \\
cell & 0.05 & Gaussian & 0.764 & 15.0 & 53 & 0.424 & 1.07 \\
fibres & 0.05 & \textbf{DCT} & 0.856 & 24.7 & 105 & 0.00862 & 0.595 \\
fibres & 0.05 & TV Bregman & 0.887 & 24.6 & 95 & 0.336 & 0.6 \\
fibres & 0.05 & Gaussian & 0.953 & 24.3 & 84 & 0.25 & 0.632 \\
vesicles & 0.05 & \textbf{DCT} & 0.889 & 26.6 & 115 & 0 & 0.539 \\
vesicles & 0.05 & TV Bregman & 0.908 & 26.5 & 111 & 0 & 0.546 \\
vesicles & 0.05 & Gaussian & 0.967 & 26.3 & 97 & 0 & 0.584 \\
\bottomrule
\end{tabular}\end{adjustbox}\end{center}
\paragraph{deconv}
\begin{center}\begin{adjustbox}{max width=\linewidth}
\begin{tabular}{lllllll}
\toprule
structure & noise & denoiser & NMSE & PSNR & SSIM & chi2\_ratio \\
\midrule
img\_001 & 0.02 & \textbf{TV Bregman} & 0.0245 & 23.4 & 0.596 & 0.98 \\
img\_001 & 0.02 & Gaussian & 0.0415 & 21.0 & 0.372 & 0.948 \\
\bottomrule
\end{tabular}\end{adjustbox}\end{center}
}{}

\paragraph{When it works, when it fails.}
\textbf{MCMC is a genuine reconstructor on MA-TIRF: it beats the ridge start} (with
$\lambda_{rr}$ in its window and a denoiser prior), and it \emph{excels} on the well-posed 2D
deconvolution, where the MMSE mean clearly beats the ridge. This \textbf{corrects} the earlier
claim (\texttt{docs/algorithms/mcmc.md}, \texttt{limits.md}) that MCMC ``cannot beat the ridge
start on MA-TIRF'' --- that limit was synthesised from lost data against a weaker configuration
($\lambda_{rr}=s_1$, the default that under-pulls). With a good $\lambda_{rr}$ MCMC is competitive
with the MAP methods, though still typically behind the plug-and-play family. Its remaining honest
limits: the denoiser must be 3D-appropriate, and \texttt{None} (no prior) cannot beat its own
start.
